% Part: incompleteness
% Chapter: arithmetization-syntax
% Section: substitution

\documentclass[../../../include/open-logic-section]{subfiles}

\begin{document}

\olfileid[pt]{inc}{art}{sub}
\olsection{Substituição}

Lembre-se de que a substituição é a operação de substituir todas as
ocorrências livres de !!a{variable}~$u$ em !!a{formula}~$!A$ por um
termo~$t$, escrita $\Subst{!A}{t}{u}$. Quando efetuada sobre
números de G\"odel de !!{variable}s, !!{formula}s e termos, essa
operação é recursiva primitiva.

\begin{prop}\ollabel{prop:subst-primrec}
Existe uma função recursiva primitiva $\fn{Subst}(x, y, z)$ com a
propriedade de que
\[
\fn{Subst}(\Gn{!A}, \Gn{t}, \Gn{u}) = \Gn{\Subst{!A}{t}{u}}.
\]
\end{prop}

\begin{proof}
Podemos então definir uma função $\fn{hSubst}$ por recursão primitiva
da seguinte maneira:
\begin{multline*}
\begin{aligned}
\fn{hSubst}(x, y, z, 0) & = \emptyseq \\
\fn{hSubst}(x, y, z, i+1) & =
\end{aligned}\\
\begin{cases}
\fn{hSubst}(x, y, z, i) \concat y & \text{se $\fn{FreeOcc}(x, z, i)$} \\
\fn{append}(\fn{hSubst}(x, y, z, i), (x)_{i}) & \text{caso contrário.}
\end{cases}
\end{multline*}
$\fn{Subst}(x, y, z)$ pode agora ser definida como
$\fn{hSubst}(x, y, z, \len{x})$.
\end{proof}

\begin{prop}
\ollabel{prop:free-for}
A relação $\fn{FreeFor}(x, y, z)$, que vale sse o termo com número de
G\"odel~$y$ é !!{free for} a variável com número de G\"odel~$z$ na
fórmula com número de G\"odel~$x$, é recursiva primitiva.
\end{prop}

\begin{proof} Exercício. \end{proof}

\begin{prob}
Prove \olref[inc][art][sub]{prop:free-for}.
\end{prob}

\end{document}
